% TOP_PROBLEM: {"id": 2800605, "title": "10 Lectures and 42 Open Problems — Constructive Kadison-Singer", "category_id": 15, "category": {"id": 15, "name": "computer_science", "display_name": "Computer Science", "description": "Computational complexity, algorithms, and theoretical CS.", "slug": "computer-science", "order_index": 15, "created_at": "2026-07-31T15:26:25.671Z"}, "status": "partially_solved"}
% FIRST_POSTED: null
% ENTRY_KIND: statement_only
% Imported from: batch08.tex; UnsolvedMath ID 2800605
% Compile twice: lualatex 2800605.tex
\documentclass[11pt]{article}
\usepackage{fontspec}
\setmainfont{Latin Modern Roman}
\usepackage{amsmath,amssymb,mathtools,unicode-math}
\setmathfont{Latin Modern Math}
\usepackage{xurl,hyperref}
\hypersetup{hidelinks,pdftitle={UnsolvedMath Problem 2800605}}
\newcommand{\sref}[2]{\hyperlink{source:#1:#2}{[S#2]}}
\newcommand{\cataloguescope}[1]{\paragraph{Mathematical question.}#1}
\begin{document}
% UnsolvedMath ID 2800605; source catalogue code AMR-027-0605
\setcounter{section}{2800604}
\section{10 Lectures and 42 Open Problems — Constructive Kadison-Singer}\label{2800605}
{\small\sffamily UnsolvedMath ID 2800605 · Computer Science}
\subsection{Definitions and mathematical statement}

\paragraph{Shared notation.}$\mathbb N=\{1,2,\ldots\}$, and $\mathbb Z,\mathbb Q,\mathbb R,\mathbb C$ denote the integers, rationals, reals and complex numbers. Any other conventions are specified locally.


A tight frame at level \(\eta\) is a finite family \(v_1,\ldots,v_m\in\mathbb C^d\) satisfying
\[\sum_{i=1}^m v_i v_i^*=\eta I_d,\]
equivalently \(\sum_i|\langle u,v_i\rangle|^2=\eta\) for every unit vector \(u\in\mathbb C^d\). Here \(v_i^*\) denotes conjugate transpose and \(\|\cdot\|_2\) is the Euclidean norm. Weaver's two-partition conclusion asks for a partition \([m]=S_1\sqcup S_2\) satisfying
\[\sum_{i\in S_j}|\langle u,v_i\rangle|^2\le\eta-\theta
\quad\text{for every unit }u\text{ and }j=1,2,\]
with universal constants \(\eta\ge2\) and \(\theta>0\), when \(\|v_i\|_2\le1\). Marcus, Spielman and Srivastava established the existence of such partitions; for example their result allows \(\eta=18\) and \(\theta=2\).
\paragraph{Statement recorded in the supplied source.}
Give a polynomial-time construction of a partition with Weaver's two-partition property: find universal constants \(\eta\ge2\), \(\theta>0\) and an algorithm that, for every finite tight frame as above with \(\|v_i\|_2\le1\), outputs \(S_1,S_2\) satisfying the displayed inequality in time polynomial in the size of its input. In particular, can the partitions guaranteed by the Kadison--Singer existence theorem be found efficiently? \sref{2800605}{2}
\subsection{Short English statement}
Construct efficiently the two-partition of any bounded tight frame guaranteed by Weaver’s formulation of the Kadison–Singer theorem, using universal constants.
\subsection{Sources}
\begin{description}
\item[\hypertarget{source:2800605:1}{S1}] ULAM.AI and the UnsolvedMath Contributors, UnsolvedMath: A Curated Collection of Open Mathematics Problems, record 2800605 (AMR-027-0605). CC BY 4.0. \url{https://huggingface.co/datasets/ulamai/UnsolvedMath}.
\item[\hypertarget{source:2800605:2}{S2}] Ben Jourdan, Peter Macgregor and He Sun, Is the Algorithmic Kadison-Singer Problem Hard? (2022), arXiv:2205.02161, Section 1, Conjecture 1 and the algorithmic question immediately following it. \url{https://arxiv.org/pdf/2205.02161}.
\end{description}
\label{end:2800605}
\end{document}
